\documentclass[../../../include/open-logic-section]{subfiles}

\begin{document}

\olfileid[hi]{sth}{ord-arithmetic}{expo}
\olsection{क्रमसूचक घातांक}

अब हम क्रमसूचक घातांक की ओर बढ़ते हैं। दुर्भाग्य से क्रमसूचक घातांक की कोई
\emph{सुंदर} संश्लेषणात्मक परिभाषा नहीं है।

निश्चय ही स्पष्ट संश्लेषणात्मक परिभाषाएँ \emph{हैं}। एक परिभाषा यह है।
$\text{finfun}(\alpha,\beta)$ उन सभी फलनों $f \colon \alpha \to \beta$
का समुच्चय हो जिनके लिए $\Setabs{\gamma \in
\alpha}{f(\gamma) \neq 0}$ किसी प्राकृतिक संख्या के साथ समसंख्य हो।
$\text{finfun}(\alpha,\beta)$ पर सुव्यवस्था इस प्रकार परिभाषित करें: $f
\sqsubset g$ तभी जब $f \neq g$ और $f(\gamma_0) < g(\gamma_0)$, जहाँ
$\gamma_0 = \text{max}\Setabs{\gamma \in \alpha}{f(\gamma) \neq
g(\gamma)}$। तब हम $\ordexpo{\alpha}{\beta}$ को
$\ordtype{\text{finfun}(\alpha, \beta), \sqsubset}$ के रूप में परिभाषित
कर सकते हैं। पॉटर इस स्पष्ट परिभाषा का उपयोग करते हैं और तुरंत समझाते हैं:
\begin{quote}
	इस क्रम का चयन केवल ऐसी क्रमसूचक घातांक की परिभाषा पाने की हमारी इच्छा
	से निर्धारित होता है जो उपयुक्त पुनरावर्ती शर्त का पालन करे\ldots;
	और इसकी कल्पना क्रमित योग या क्रमित गुणनफल की अपेक्षा कहीं कठिन है।
	\citep[p.~199]{Potter2004}
\end{quote}
ठीक यही बात है। हमने योग को ``$\alpha$ की एक प्रति जिसके बाद $\beta$ की
एक प्रति हो'' और गुणन को ``$\alpha$ की प्रतियों का एक $\beta$-अनुक्रम''
बताया था। किंतु $\text{finfun}(\alpha, \gamma)$ के बारे में कहने के लिए
हमारे पास ऐसा कोई संक्षिप्त वर्णन नहीं है। इसलिए इसके बदले हम क्रमसूचक
घातांक की परिभाषा सीधे परिमितेतर पुनरावर्तन \emph{द्वारा} देंगे, अर्थात्:

\begin{defn}\ollabel{ordexporecursion}
\begin{align*}
	\ordexpo{\alpha}{0} &= 1\\
	\ordexpo{\alpha}{\beta\ordplus 1} &=\ordexpo{\alpha}{\beta} \ordtimes \alpha\\
	\ordexpo{\alpha}{\beta} &= \bigcup_{\delta < \beta}\ordexpo{\alpha}{\delta}& & \text{जब $\beta$ सीमा क्रमसूचक हो}
\end{align*}
\end{defn}

यदि हम समुच्चय सिद्धांतकारों \emph{के रूप में} काम कर रहे होते, तो शायद
क्रमसूचक घातांक के कुछ गुणों का अन्वेषण करना चाहते। किंतु इसके अतिरिक्त
कहने को अधिक कुछ नहीं है, सिवाय इस अपेक्षित तथ्य के कि क्रमसूचक घातांक
क्रमविनिमेय नहीं है। अतः $\ordexpo{2}{\omega} =
\bigcup_{\delta < \omega}\ordexpo{2}{\delta} = \omega$, जबकि
$\ordexpo{\omega}{2} = \omega \ordtimes \omega$। पर हमें घातांक के
क्रमविनिमेय होने की \emph{अपेक्षा} भी नहीं करनी चाहिए, क्योंकि प्राकृतिक
संख्याओं के लिए भी वह क्रमविनिमेय नहीं है:
$\ordexpo{2}{3} = 8 < 9 = \ordexpo{3}{2}$।

\begin{prob}
परिमितेतर आगमन का उपयोग करके सिद्ध कीजिए कि यदि हम
$\ordexpo{\alpha}{\beta} = \ordtype{\text{finfun}(\alpha, \beta),
\sqsubset}$ परिभाषित करें, तो हमें
\olref[sth][ord-arithmetic][expo]{ordexporecursion} के पुनरावर्तन समीकरण
प्राप्त होते हैं।
\end{prob}

\end{document}
